\section{Limitations}
\label{sec:limitations}

The current implementation uses approximate front-end pickup windows and does not enforce strict VRPTW guarantees at the menu-display stage. Final route costs are recovered through HGS-based re-optimization \citep{vidal2022hybrid}, which makes the system suitable for comparative policy analysis but not for claims about exact optimal routing. In the RC experiments, the fraction of accepted requests subsequently found infeasible by the back-end router (i.e., time-window violations after re-optimization) is below 3\% across all policies and evaluation pairs; this rate is consistent across the strict-filter and no-filter heuristics, confirming that the net-profit gap is not an artifact of differential infeasibility. Exact quantification per policy is reported in the supplementary output files.

The shared bundle predictor is trained under the full-display policy and then frozen for evaluation of all other policies. This introduces a potential distribution shift: the menu distribution induced by the no-filter heuristic differs from the training distribution, and the predictor's cost and ETA estimates may carry systematic bias when applied to bundles selected under a different filtering regime. The current logging partially exposes this shift. In the current RC outside-option study, the mean predicted cost of the chosen bundle is $43.53$ under full display, $43.42$ under the strict-filter heuristic, and $43.36$ under the no-filter heuristic. We also record displayed-offer ETA/IVT error: displayed ETA MAE is $2577$\,s under full display, $2564$\,s under the strict-filter heuristic, and $2384$\,s under the no-filter heuristic; the corresponding IVT MAE values are $2308$\,s, $2284$\,s, and $2131$\,s. These diagnostics are descriptive rather than causal, but they confirm that removing the filter changes the evaluated bundle distribution and should be interpreted jointly with the policy outcomes. The quantile error diagnostics (Table~\ref{tab:filter_validity_summary}) and false-negative pruning breakdown by distance band now provide a more complete picture of filtering behavior: ETA P95 errors of approximately 3\,600 seconds indicate a coarse filtering signal, and the concentration of false-negative pruning in the far distance band ($>$1\,500\,m) confirms that filtering is operationally conservative for nearby meeting points. However, a true filtering confusion matrix requiring ground-truth feasibility labels remains unavailable, so the filter-validity evidence is still partially indirect. The large prediction errors documented above represent a key limitation: while the FN concentration pattern suggests that filtering is operationally conservative for nearby meeting points, the coarse prediction signal means that the operational reliability of filtering-based assortment decisions cannot be fully established from the current diagnostics alone.

The empirical evaluation covers three instance types: the RC synthetic benchmark (Homberger--Gehring family) and the Austin and Seattle real-world-derived instances adapted from the Amazon Last Mile Routing Research Challenge \citep{merchan2022amazon}, each with two bundled splits. While this provides cross-topology evidence spanning synthetic and real-world-derived geographies, all instances share the same fleet size (15 vehicles, capacity 25) and the same demand intensity regime. The city studies are also low-sample: only two split pairs are available per city, so the Austin and Seattle results should be interpreted as suggestive external checks rather than definitive replication. The operational-baseline comparison in Table~\ref{tab:operational_baselines} covers only the RC low-uptake regime and does not span the calibrated higher-uptake or city instances; extending this reference to additional demand conditions is left as future work. Broader coverage --- including larger fleets, additional real-world-derived instances, and comparison with published DRT assignment methods --- is left as explicit future work.

The MNL utility parameters ($\beta$, $\bar{u}$, $u^{\text{home}}$, $\alpha_\tau$, $\alpha_\Delta$, $\sigma$) are fixed design parameters of the simulator rather than values estimated from passenger survey data. This matters directly for interpretation. Under the original RC outside-option evaluation, the benchmark exhibits extremely low realized uptake of non-home bundles; under the calibrated higher-uptake regime, the same simulator structure yields materially higher acceptance. The paper therefore treats calibration as an evidence-setting decision rather than as a resolved demand-estimation result. The higher-uptake regime is useful because it restores behavioral variation, but it should not be read as a uniquely validated empirical demand model. The outside-option utility scan (Section~\ref{sec:outside_option_sensitivity}) narrows this limitation by demonstrating that the results are not obviously brittle within this RC stress test across five $u_0$ levels, but it does not eliminate the fundamental demand-calibration gap. Appendix Table~\ref{tab:mnl_sensitivity_summary} records the current sensitivity evidence and the reproducible suite for filling the remaining low-, medium-, and higher-uptake cells under the current pipeline.

Consumer surplus is computed under the implemented MNL utility specification (Appendix~\ref{app:consumer_surplus}). The all-request surplus normalizes outside-option surplus to zero, while the accepted-user surplus averages only over requests where a menu bundle was selected. When comparing policies, a change in acceptance rate shifts the composition of accepted users: if a filter removes low-value requests, the accepted-user surplus can increase even if no individual passenger's experience improves. This selection effect is inherent in the menu-optimization framework and should be considered when interpreting surplus differences across policies.

The seven scoring and menu parameters ($\mu_m, \mu_w, \mu_t, \mu_v, \mu_r, \mu_q, k$) are similarly fixed design choices. Exact enumeration is exact only for the implemented surrogate objective on bounded candidate sets; it is not an exact solution to the full dynamic DRT problem with downstream HGS re-optimization. The pricing-baseline comparison is informative about one calibrated stress-test regime, but it is still a local result: flat markdown can outperform the Lambert-W-inspired rule in the reported higher-uptake environment, yet we do not claim that this ordering is invariant to other revenue targets, price caps, demand scales, or preference settings. The outside-option scan (Section~\ref{sec:outside_option_sensitivity}) similarly demonstrates directional consistency within the tested $u_0$ range but does not establish robustness to all possible demand specifications. A single-split sensitivity table from an earlier pipeline run is retained in Appendix~\ref{app:sensitivity} as an archival diagnostic, but it was not rerun under the current evaluation setup and should not be interpreted as confirmatory evidence that one solver, filter, or pricing rule robustly dominates the alternatives.
